\scanpage{24}
Considérons la partition $\Pi_I^{\unsure{\circ|0}}$ de $X\setminus\bigcup_{\alpha'\in I}E_{\alpha'}$, soit $S$ l'ens des classes, notées $U_s$ $(s\in S)$. \deleted{\unsure{On a|On}} Considérons la \deleted{\uline{Lemme} $\forall i\in I$,} \inserted{Relation $R$ entre $S$ et $I$}
\[
\deleted{\unsure{U|E}}(s,i)\in R\ \overset{\text{def}}{\Longleftrightarrow}\ \overline{U_s}\cap E_i\ne\varnothing.
\]

\uline{Lemme} Cette relation fait de $(S,I\unsure{)\deleted{R})|\deleted{)}{,}R)}$

\noindent\begin{tabular}{@{}|p{.94\linewidth}@{}}
un graphe strict ayant $S$ comme \unsure{ens des|ens de} sommets, $I$ comme ens. \unsure{de|d'} \unsure{arêtes|arètes|aretes}.) \deleted{Cette relation} \unsure{Le|Ce} graphe est un \uline{arbre}
\end{tabular}

le fait qu'on a un graphe strict signifie que a) $\forall i\in I$, Il \deleted{\unsure{existe|ex}} y a exactement deux $U_s$ tels que $\overline{U_s}\cap E_i\ne\varnothing$, et b) $i$ est caractérisé par l'ens de ces deux $U_s$.

\unsure{Soient|Soit} $\Sigma_i^+,\Sigma_i^-$ les deux comp. \unsure{connexes} de $X\setminus E_i$, alors a) se précise en : il y a \unsure{exactement|exactemt|existe} \uline{un} $U_{s_{\unsure{+|-}}}$ contenu dans $\Sigma_i^{\unsure{+|-}}$ tel que $\overline{U_{s_{\unsure{+|-}}}}\cap E_i\ne\varnothing$ (et de $=$ pour $\Sigma_i^-$, \unsure{dé|qui} \unsure{définit|definit} un $U_{s_-}\subset\Sigma_i^-$ tel que $\overline{U_{s_-}}\cap E_i\ne\varnothing$),
% archive sheet number: 23
